\chapter{Math}

\section{Choose the reduction}
Huge integer + allowed error: bound discarded digits before factoring.
Many occurrences, values $\le V$: histogram/multiples, not all pairs.
Decimal division: reduce the denominator, then remainder states.
Sum over subarrays/pairs: count each feature's contribution.
Repeated operations: seek an invariant; necessity is not sufficiency.
Random process: indicators or expectation recurrence before simulation.

\section{Exact arithmetic}
\kactlimport{Prelude.h}
\subsection{Floor, ceiling and bounds}
C++ division truncates toward zero: $-7/3=-2$, but
$\lfloor-7/3\rfloor=-3$. For $a\ge0,b>0$,
$\lceil a/b\rceil=a/b+[a\bmod b\ne0]$; avoids overflow in $a+b-1$.
For positive $d$, multiples in $[L,R]$ correspond to integer indices
$[\lceil L/d\rceil,\lfloor R/d\rfloor]$; answer is zero if empty.

Widen \emph{before} arithmetic: \texttt{i128(a)*b},
\texttt{i128(b.x)-a.x}. Assignment to a wide result is too late.
Signed 64-bit maximum is about $9.22\cdot10^{18}$.
Bound the whole sum, not only a term: pair counts $O(n^2)$,
all subarray sums $O(n^3 A)$ for $|a_i|\le A$.
Helpers reject division by zero and an unrepresentable quotient.
\texttt{i128} requires GCC/Clang; not arbitrary precision.
\kactlimport{IntegerDivision.h}

\subsection{Exact decimal input}
For at most $k$ fractional places, store the text scaled by $10^k$.
Cents: \texttt{12.05}$\to1205$, \texttt{3}$\to300$,
\texttt{-0.07}$\to-7$. Read a string, not a \texttt{double}.
Parser rejects excess fractional places; it does not round.
No exponent notation or commas. Digit count excludes the sign;
zero has one digit. Integer division by ten is exact at powers of ten;
\texttt{log10} can misclassify a boundary.
\kactlimport{ExactDecimals.h}

\section{Congruences}
\subsection{Powers and inverses}
Normalize residues into $[0,m)$, $m>0$.
Addition/multiplication preserve congruence; division usually does not.
An inverse of $a$ exists iff $\gcd(a,m)=1$.
For prime $p$ and $p\nmid a$, $a^{-1}=a^{p-2}\pmod p$.
Composite modulus: use extended Euclid.
Example $3^{-1}=3\pmod8$, but $2^{-1}\pmod8$ does not exist.
Binary powering: $O(\log e)$ time, $O(1)$ space, $e\ge0$.
A widened product is needed even when both residues fit 64 bits.
Our exponent-zero convention gives $0^0=1$ before reduction.
\kactlimport{ModularArithmetic.h}

\subsection{Bezout and shared-factor CRT}
Extended Euclid gives $ax+by=g=\gcd(a,b)$ in $O(\log\min(a,b))$.
$ax+by=c$ has integer solutions iff $g\mid c$.
From one solution $(x_0,y_0)$, all are
$(x_0+t b/g,\ y_0-t a/g)$, $t\in\mathbb Z$.
Use floor/ceiling to intersect bounds on $t$.
If $a=b=0$, handle $c=0$ separately.

Merge $X\equiv a\pmod m$, $X\equiv b\pmod n$:
$g=\gcd(m,n)$ must divide $b-a$. Put $X=a+mt$ and solve
\[
 (m/g)t\equiv(b-a)/g\pmod{n/g}.
\]
Answer is unique modulo $\operatorname{lcm}(m,n)$.
Example $(a,m)=(2,6),(b,n)=(5,9)$ gives $X=14\pmod{18}$;
residue $4$ modulo $9$ is incompatible.
Merge many equations one by one. Positive moduli; modulus one adds
no restriction. Helper throws if the merged lcm exceeds 64 bits.
\kactlimport{EuclidCRT.h}

\section{Prime exponents and counting values}
\subsection{Factorization and totient}
For $n=\prod p^{e_p}$,
\[
 \tau(n)=\prod(e_p+1),\quad
 \sigma(n)=\prod(1+p+\cdots+p^{e_p}),
\]
\[
 \phi(n)=n\prod_{p\mid n}(1-1/p).
\]
Example $72=2^3 3^2$: $\tau=12$, $\sigma=195$, $\phi=24$.
For one, the empty products give all three values one.
Generate divisors by extending old divisors with each power of a new prime.
Geometric accumulation avoids division modulo a composite modulus.

Trial division stops at $p\le n/p$, reducing $n$ as factors are removed;
remaining $n>1$ is prime. Worst case $O(\sqrt n)$.
A ten-digit core needs candidates only through $10^5$;
repeated arbitrary 64-bit factorization needs a different algorithm.
SPF sieve: $O(V)$ time/space; factor $x\le V$ in $O(\log x)$.
Zero/negative inputs require separate definitions.
\kactlimport{FactorSieve.h}

\subsection{Histogram and exact gcd classes}
For positive values $\le V$, let $f[x]$ count occurrences and
$M[d]=\sum_{k\ge1}f[kd]$. All $M$ take $O(V\log V)$ time,
$O(V)$ memory. Zero needs separate handling.
For fixed $b$ and divisors $g\mid b$, process decreasing $g$:
\[
 E_b[g]=M[g]-\!\sum_{\substack{h\mid b,\ g\mid h\\h>g}} E_b[h].
\]
Then $E_b[g]$ counts numerators with $\gcd(a,b)=g$ exactly.
For values $[2,3,6,6]$, $b=6$: $E[6]=2,E[3]=1,E[2]=1,E[1]=0$.
For ordered occurrence pairs multiply by $f[b]$.
Nested divisor work needs its own bound, e.g. $\sum_b\tau(b)^2$;
it is not automatically $O(V\log V)$.

\section{Remainder states and decimals}
\subsection{Digits as a small state}
Append digit $d$ to remainder $r$ modulo $m$:
$r'=(10r+d)\bmod m$. Widen $10r$ if needed.
At the same position, equal remainder means the same possible suffixes
when no other constraint distinguishes prefixes. Hence a digit-pattern DP
uses $(\mathrm{position},r)$, with $O(10Lm)$ time for length $L$.
Leading-zero rules and extra digit restrictions must enter transitions/state.
Smallest reconstruction: suffix feasibility, then try digits in increasing
order while retaining a feasible suffix. Counting and existence differ.

\subsection{Period and rational display}
For $\gcd(a,m)=1$, multiplicative order $\operatorname{ord}_m(a)$
is the smallest $k>0$ with $a^k\equiv1\pmod m$.
It divides $\phi(m)$; equality need not hold.
Start $k=\phi(m)$; for prime $p\mid k$, repeatedly divide by $p$
when $a^{k/p}\equiv1\pmod m$. Factorization dominates large inputs.
Example $\operatorname{ord}_9(10)=1$, while $\phi(9)=6$.

For $a\ge0,b>0$, reduce $q=b/\gcd(a,b)=2^u5^v r$,
$\gcd(r,10)=1$. Canonical displayed digit count (no point or overbar):
\[
 D(a,b)=\mathrm{digits}(\lfloor a/b\rfloor)+\max(u,v)+P(r),
\]
where $P(1)=0$, else $P(r)=\operatorname{ord}_r(10)$.
Reason: long division advances remainder by multiplication by ten;
after the $2,5$ prefix, a repeated remainder requires $10^k\equiv1$.
Do not count artificial trailing zeros or the alternative recurring-nine form.

Examples: $1/4$: $1+2+0=3$; $1/6$: $1+1+1=3$;
$3/14$: $1+1+6=8$; $41/4$: $2+2+0=4$.
Zero has one integer digit.
\emph{Displaying Decimals (2025 D):} separate fractional length by gcd
classes; integer digits add one per threshold $a\ge10b,100b,\ldots$,
answered with frequency suffix counts.
For $[15,36,14]$, nine ordered lengths sum to $30$; expectation $10/3$.
The single-fraction formula passed direct long-division checks;
full maximum-size aggregation is not benchmarked here.
\kactlimport{DecimalPeriod.h}

\section{Approximation as a construction}
\subsection{Fuzzy Factorization (2025 F)}
Allowed relative error $10^{-9}$ changes the task. If $X$ has more than ten
digits, keep ten leading digits $m$ and discard $e$ places:
\[
 X=m10^e+r,\quad 0\le r<10^e,\quad m\ge10^9.
\]
Choose $Y=m10^e$. Then
\[
 \frac{|X-Y|}{X}<\frac1m\le10^{-9}.
\]
Factor only $m<10^{10}$; add $e$ to exponents of $2$ and $5$.
Merge repeated primes. For at most ten digits choose $Y=X$.
Time $O(|X|+\sqrt m)$; no huge integer arithmetic.

Example $X=123456789012345$: $m=1234567890$, $e=5$,
$m=2\cdot3^2\cdot5\cdot3607\cdot3803$.
Thus $Y=2^6 3^2 5^6\cdot3607\cdot3803$.
Discarded remainder $12345$ obeys $12345\cdot10^9\le X$.
Tolerance $10^{-k}$ similarly permits $k+1$ leading digits, but verify
factorization cost and output prime limits. Empty factor map means one;
adapt the required output format. Positive input, no leading zeros.
\kactlimport{FuzzyFactorization.h}

\section{Combinatorics}
\subsection{Choose the sample objects}
$A$ choices each followed by exactly $B$ choices give $AB$.
If counts vary, sum them. Disjoint cases add; overlaps need correction.
Ordered distinct $k$-tuples: $n!/(n-k)!$; unordered subsets:
$\binom nk=n!/(k!(n-k)!)$; ordered repetition: $n^k$.
Five people: president/secretary $20$, two-person committee $10$.
Stars and bars: $k$ ordered bins, total $n$ identical items;
nonempty $\binom{n-1}{k-1}$, empty allowed $\binom{n+k-1}{k-1}$.
For $k\ge1$; upper capacities need additional reasoning.

Pascal: $\binom nk=\binom{n-1}{k-1}+\binom{n-1}k$,
edge values one. Works modulo any positive modulus: $O(n^2)$ time,
$O(n)$ memory; update a single row in decreasing $k$.
Prime factorial tables: $O(N+\log p)$ preprocessing, $O(1)$ queries,
$O(N)$ memory. Require prime $p$ and $N<p$.
At $N\ge p$, $N!$ is zero modulo $p$; composite moduli may fail earlier.
The constructor trusts primality. No automatic Lucas/composite support.
\kactlimport{Combinations.h}

\subsection{Contributions and complements}
Zero-based index $i$ appears in $(i+1)(n-i)$ subarrays, so
\[
 \sum_{\mathrm{subarrays}}\sum a_i
 =\sum_i a_i(i+1)(n-i).
\]
Example $[2,1,3]$: $6+4+9=19$. Widen all products.
For unordered pairs:
\[
 \sum_{i<j}(a_i+a_j)=(n-1)\sum_i a_i,
\]
\[
 \sum_{i<j}(a_i-a_j)^2=n\sum_i a_i^2-(\sum_i a_i)^2.
\]
At least one good feature = total minus zero good features.
A maximal bad run of length $L$ contributes $L(L+1)/2$ all-bad
substrings. Circular input needs merged boundary runs and a separate
all-bad case; the linear formula alone is insufficient.

\subsection{Inversions and repeated copies}
Inversion: $i<j$, $a_i>a_j$ strictly. Equal values do not count.
Fenwick scan adds previously seen minus previous values $\le a_i$.
Coordinate compression preserves comparisons; $O(n\log n)$ time/space $O(n)$.
For $K$ copies of $S$:
\[
 \mathrm{inv}(S^K)=K\,\mathrm{within}+\binom K2\,\mathrm{cross},
\]
where $\mathrm{cross}=\#\{(u,v):S_u>S_v\}$ without an index-order
restriction inside one copy. Every earlier-copy position precedes every
later-copy position. For $S=ba,K=3$, answer $3+3=6$.
This is the decomposition in 2023 I. Use wide/modular products.
\kactlimport{Inversions.h}

\section{Probability and expectation}
Uniform outcomes must be specified: ordered occurrence pairs with
replacement number $N^2$, without replacement $N(N-1)$.
Equal values may still be different occurrences.
\[
 \mathbb E[X]=\sum_x x\Pr(X=x),\qquad
 \mathbb E[\sum I_j]=\sum\Pr(I_j=1).
\]
Linearity needs no independence. Multiplying probabilities does:
$\Pr(A\cap B)=\Pr(A)\Pr(B\mid A)$.
Random permutation expected inversions: each pair has probability $1/2$,
so $n(n-1)/4$. Expectation need not be a possible outcome.

Self loop: $E=c+qE+\sum_jp_jE_j$ gives
$E=(c+\sum_jp_jE_j)/(1-q)$ for $q<1$.
A permanent positive-cost self loop has infinite expectation.
Modulo prime $p$, fraction $A/B$ means $A B^{-1}$ only if $p\nmid B$;
reduce exact cancellations first if needed. A real finite expectation need
not have an invertible modular denominator. Never modularize a decimal
approximation. $\mathbb E[X^2]\ne\mathbb E[X]^2$ generally;
$\operatorname{Var}(X)=\mathbb E[X^2]-\mathbb E[X]^2$.

\section{Proof and inequality checks}
An invariant needs initial truth and preservation by every operation.
Try sum, parity, weighted sum, gcd, multiset, order.
Invariant equality is necessary; prove sufficiency by a construction.
A decreasing nonnegative integer proves termination.
Binary search needs a monotone predicate, not just a fast checker.
Greedy needs an exchange preserving feasibility and not worsening cost.
Brute-force tiny cases find counterexamples; passing them is not proof.

Useful sums:
$\sum_{i=1}^n i=n(n+1)/2$,
$\sum_{i=0}^k r^i=(r^{k+1}-1)/(r-1)$ for $r\ne1$.
Before modular division, cancel an exact factor or verify invertibility.
For error bounds compare with the original denominator; check strict vs
nonstrict inequalities and small inputs separately.

\section{Integer geometry}
For vectors $u,v$,
$u\cdot v=u_xv_x+u_yv_y$,
$u\times v=u_xv_y-u_yv_x$.
Dot zero: perpendicular (nonzero vectors); positive/negative: acute/obtuse.
$\operatorname{orient}(A,B,C)=(B-A)\times(C-A)$:
positive left, negative right, zero collinear.
Compare squared distances; no square root needed.
Twice triangle area is $|\operatorname{orient}|$.

Point on segment: orientation zero and both coordinates within inclusive
endpoint ranges. Degenerate $A=B$ works. Segment intersection: opposite
orientation signs for both lines, plus on-segment zero cases; collinear
segments reduce to interval overlap. Compare signs, do not multiply cross
products. Polygon in boundary order:
$2A=|\sum_i P_i\times P_{i+1}|$ for a simple polygon.
Self intersections give algebraic area; sorting by $x$ is not boundary order.

Helpers use 128-bit operations and require coordinate magnitudes
$\le10^{18}$. Arbitrary full-range 64-bit coordinates can overflow even
128 bits; polygon accumulation needs a separate bound in $n$.
Integer predicates are exact; real intersections/distances need tolerances
based on the statement, not a universal EPS.
\kactlimport{IntegerGeometry.h}
